\documentclass[11pt]{article}
\usepackage[margin=1in]{geometry}
\usepackage{enumitem}
% =============================================================================
% Preambles/header.tex — shared LaTeX/Beamer preamble for this project
%
% Use from a Beamer lecture by prepending:
%     \input{header}
% and compiling with TEXINPUTS=../Preambles:$TEXINPUTS (the /compile-latex
% skill does this automatically).
%
% PALETTE CONTRACT. Color names defined here MUST match the names in
% Quarto/theme-template.scss so Beamer and Quarto renderings use the same
% palette. The scripts/check-palette-sync.sh script enforces this.
%
% To customize: edit the HEX values below AND the matching $variable in
% Quarto/theme-template.scss, then run scripts/check-palette-sync.sh.
% =============================================================================

% -----------------------------------------------------------------------------
% Engine detection + encoding
%
% Our /compile-latex skill uses XeLaTeX, where inputenc is unnecessary and
% can produce encoding warnings. Only load inputenc under pdfTeX.
% -----------------------------------------------------------------------------
\usepackage{iftex}
\ifPDFTeX
  \usepackage[utf8]{inputenc}
\fi

\usepackage{amsmath, amssymb, amsthm}
\usepackage{booktabs}
\usepackage{graphicx}

% -----------------------------------------------------------------------------
% PALETTE — mirrors Quarto/theme-template.scss variable names.
% Load xcolor and define colors BEFORE anything that references them
% (notably hyperref, hypersetup, and the Beamer color assignments).
% -----------------------------------------------------------------------------
\usepackage{xcolor}

\definecolor{primary-blue}{HTML}{012169}      % headings, accents
\definecolor{primary-gold}{HTML}{B9975B}      % emphasis, borders
\definecolor{highlight-yellow}{HTML}{F2A900}  % markers, alerts
\definecolor{light-bg}{HTML}{E8EDF5}          % title / section backgrounds
\definecolor{jet}{HTML}{1A1A1A}               % body text

% Semantic colors (used in diagrams and annotations)
\definecolor{positive}{HTML}{15803D}          % good / observed / identified
\definecolor{negative}{HTML}{B91C1C}          % bad / problematic / bias
\definecolor{neutral}{HTML}{525252}           % reference / context

% Highlight accents (match .hi-* classes in the SCSS)
\definecolor{hi-slate}{HTML}{314F4F}
\definecolor{hi-green}{HTML}{2E8B57}
\definecolor{hi-red}{HTML}{C41E3A}

% Semantic aliases so skill templates and snippets can write
% \textcolor{accent}{...} without hard-coding "primary-blue".
\colorlet{accent}{primary-blue}
\colorlet{accent2}{primary-gold}

% -----------------------------------------------------------------------------
% Hyperref — loaded AFTER xcolor + palette so link colors resolve.
% -----------------------------------------------------------------------------
\usepackage{hyperref}
\hypersetup{colorlinks=true, linkcolor=primary-blue, urlcolor=primary-blue,
            citecolor=primary-blue, pdfborder={0 0 0}}

% -----------------------------------------------------------------------------
% Course code — set once per deck (after \input{header}, before \begin{document})
% via \coursecode{CS401}. Shown in the footer on every slide so a deck's course
% is identifiable without opening the title page. Empty by default.
% -----------------------------------------------------------------------------
\makeatletter
\newcommand{\coursecode}[1]{\gdef\@coursecodetext{#1}}
\gdef\@coursecodetext{}
\makeatother

% -----------------------------------------------------------------------------
% Beamer theme assignments (only apply under Beamer).
%
% We detect Beamer via \@ifclassloaded{beamer}, which is the documented,
% non-internal way. This must be wrapped in \makeatletter/\makeatother
% because \@ifclassloaded contains an @-catcode character.
% -----------------------------------------------------------------------------
\makeatletter
\@ifclassloaded{beamer}{%
  \usefonttheme{professionalfonts}
  \setbeamercolor{structure}{fg=primary-blue}
  \setbeamercolor{frametitle}{fg=primary-blue}
  \setbeamercolor{title}{fg=primary-blue}
  \setbeamercolor{subtitle}{fg=primary-gold}
  \setbeamercolor{author}{fg=primary-blue}
  \setbeamercolor{institute}{fg=neutral}
  \setbeamercolor{date}{fg=primary-gold}
  \setbeamercolor{itemize item}{fg=primary-blue}
  \setbeamercolor{itemize subitem}{fg=primary-gold}
  \setbeamercolor{alerted text}{fg=primary-gold}
  \setbeamercolor{block title}{fg=white, bg=primary-blue}
  \setbeamercolor{block body}{bg=light-bg}
  % Semantic block variants: block=definition/concept (blue), exampleblock=
  % worked example (green/positive), alertblock=key takeaway or caution (gold).
  % Keep to <=2 colored boxes per slide across ALL three variants (INV-7).
  \setbeamercolor{block title example}{fg=white, bg=positive}
  \setbeamercolor{block body example}{bg=light-bg}
  \setbeamercolor{block title alerted}{fg=white, bg=primary-gold}
  \setbeamercolor{block body alerted}{bg=light-bg}

  % Minimal footer: course code (left) + page X / N (right), no navigation clutter
  \setbeamertemplate{navigation symbols}{}
  \setbeamertemplate{footline}{%
    \color{neutral}\footnotesize\hspace{0.7em}\@coursecodetext\hfill
    \insertframenumber/\inserttotalframenumber\hspace{0.7em}\vspace{0.3em}}
}{}
\makeatother

% -----------------------------------------------------------------------------
% TikZ setup — libraries the snippet gallery uses
% -----------------------------------------------------------------------------
\usepackage{tikz}
\usetikzlibrary{arrows.meta, positioning, calc, decorations.pathreplacing,
                fit, shapes.geometric, backgrounds}

% Shared TikZ styles — snippets and hand-written diagrams can pick these up
% via `\begin{tikzpicture}[every node/.style={...}]` or by naming specific
% styles (e.g., `\node[dag-node] (X) at (0,0) {X};`).
\tikzset{
  % Node styles for causal DAGs and similar
  dag-node/.style = {
    draw=primary-blue, fill=light-bg, rounded corners=2pt,
    minimum width=1.2cm, minimum height=0.9cm, align=center, font=\small
  },
  decision-node/.style = {
    draw=primary-gold, fill=white, diamond, aspect=1.8,
    minimum width=1.8cm, minimum height=1.2cm, align=center, font=\small,
    inner sep=1pt
  },
  % Edge styles — observed vs counterfactual
  observed-edge/.style = {-{Latex[length=2.5mm]}, thick, draw=jet},
  counterfactual-edge/.style = {-{Latex[length=2.5mm]}, thick, draw=neutral, dashed},
  confound-edge/.style = {-{Latex[length=2.5mm]}, thick, draw=negative, dashed},
  % Data-point styles for plots
  observed-dot/.style = {fill=primary-blue, draw=primary-blue, circle, inner sep=1.2pt},
  counterfactual-dot/.style = {fill=white, draw=neutral, circle, inner sep=1.2pt}
}

% -----------------------------------------------------------------------------
% Convenience macros
% -----------------------------------------------------------------------------
\newcommand{\muted}[1]{\textcolor{neutral}{#1}}
\newcommand{\key}[1]{\textcolor{primary-gold}{\textbf{#1}}}
\newcommand{\good}[1]{\textcolor{positive}{#1}}
\newcommand{\bad}[1]{\textcolor{negative}{#1}}

% Standout frame macro: full-bleed dark background for section transitions.
% Beamer-only (uses \begin{frame}/\setbeamercolor/\usebeamerfont) -- do not
% call from an article-class file (e.g. Notes/<CODE>/*-notes.tex). \muted,
% \key, \good, \bad, and \coursecode above are plain xcolor/gdef and are
% safe to use from either class; verified when Notes/article-class support
% was added.
% Expand: \transitionslide{Your transition title}
\newcommand{\transitionslide}[1]{%
  \begingroup
  \setbeamercolor{background canvas}{bg=primary-blue}
  \begin{frame}[plain]
    \centering
    \vfill
    {\usebeamerfont{title}\color{white}\Large #1\par}
    \vfill
  \end{frame}
  \endgroup
}

% Section divider frame (Beamer-only), an alternative to \transitionslide with a
% darker two-line style: near-black (jet) background, a small white label line
% above a larger gold title, and a thin gold rule along the bottom edge.
% Expand: \sectiondivider{Section 2}{Pointers}
\newcommand{\sectiondivider}[2]{%
  \begingroup
  \setbeamercolor{background canvas}{bg=jet}
  \begin{frame}[plain]
    \vspace*{3.0cm}
    {\color{white}\ttfamily\large #1\par}
    \vspace{0.5cm}
    {\color{primary-gold}\LARGE #2\par}
    \begin{tikzpicture}[remember picture, overlay]
      \draw[primary-gold, line width=2.5pt]
        ($(current page.south west)+(0,0.1)$) -- ($(current page.south east)+(0,0.1)$);
    \end{tikzpicture}
  \end{frame}
  \endgroup
}

% -----------------------------------------------------------------------------
% Channel watermark — "Concepts That Click" (YouTube).
%
% Article-class files (CompetitiveExam Q/A sets, Notes, Assignments) get a
% light diagonal draft watermark on every page plus a centered footer naming
% the channel. Beamer decks get a subtle TikZ background watermark instead
% (the footline already carries the course code). Centralized here so every
% MCQ PDF is branded without per-file edits.
% -----------------------------------------------------------------------------
\makeatletter
\@ifclassloaded{article}{%
  \usepackage{draftwatermark}
  \SetWatermarkText{Concepts That Click}
  \SetWatermarkScale{1}
  \SetWatermarkFontSize{2cm}
  \SetWatermarkLightness{0.86}
  \SetWatermarkAngle{45}
  \usepackage{fancyhdr}
  \pagestyle{fancy}
  \fancyhf{}
  \fancyfoot[C]{\footnotesize\textcolor{neutral}{Concepts That Click~|~YouTube: \href{https://www.youtube.com/@concepts.thatclick}{@concepts.thatclick}}}
  \renewcommand{\headrulewidth}{0pt}
  \renewcommand{\footrulewidth}{0pt}
  \fancypagestyle{plain}{%
    \fancyhf{}%
    \fancyfoot[C]{\footnotesize\textcolor{neutral}{Concepts That Click~|~YouTube: \href{https://www.youtube.com/@concepts.thatclick}{@concepts.thatclick}}}%
    \renewcommand{\headrulewidth}{0pt}%
    \renewcommand{\footrulewidth}{0pt}%
  }
}{}
\@ifclassloaded{beamer}{%
  \setbeamertemplate{background}{%
    \begin{tikzpicture}[remember picture, overlay]
      \node[rotate=45, opacity=0.055, align=center]
        at (current page.center)
        {\Huge\textcolor{neutral}{Concepts That Click}};
    \end{tikzpicture}%
  }
}{}
\makeatother 
\coursecode{JKSSB}
\renewcommand{\thesection}{13.\arabic{section}}
\newtheorem{example}{Example}[section]
\newenvironment{definitionbox}[1]{%
  \par\noindent\textbf{Definition (#1).}\ \itshape
}{\par\vspace{0.3em}}
\title{OCR, OMR, MICR --- Detailed Lecture Notes}
\author{JKSSB Computer Awareness}
\date{\today}

\begin{document}
\maketitle

\section{Reading machines and why they matter}
A computer works only with data that has already been converted into
machine-readable form. When the data exists on paper --- as printed text, as
a filled-in bubble on a form, or as a line of characters on a cheque --- the
machine needs a way to read it. A \textbf{reading machine} is the device or
technology that captures this physical information and turns it into digital
data the system can process. The physical reader (often a scanner, sometimes
a dedicated reader) performs the capture; the \emph{recognition method}
decides what the captured pattern means. This lesson covers the three most
examined recognition methods --- OCR, OMR, and MICR --- together with the
code-and-tag devices (barcode, QR code, RFID) and biometric input, because
questions frequently ask candidates to tell these near-neighbours apart.

The single most important idea is the \emph{read-target} of each method.
\textbf{OCR} reads \emph{text}. \textbf{OMR} reads \emph{marks}.
\textbf{MICR} reads \emph{characters printed in magnetic ink}. Almost every
item on this topic can be solved by fixing this one distinction in memory.

All the devices in this lesson are \emph{input} devices or reading
\emph{technologies}: they bring data into the system. Most of them are also
\emph{direct} input, because the device reads the object, mark, code, or
trait itself, instead of asking the user to type it. This is why a reading
machine is faster and less error-prone than a keyboard for a repetitive task
such as reading a thousand answer sheets or a thousand cheques.

\begin{definitionbox}{Reading machine}
A device or technology that captures printed characters, marks, codes, or
tags from a physical medium and converts them into machine-readable data for
the computer to process.
\end{definitionbox}

\section{Terminology at a glance}
The table below fixes the meaning of every term used in this lesson. Read it
down once; then, for each term, be ready to state the read-target and one
place it is used.

\begin{center}
\begin{tabular}{p{0.24\textwidth}p{0.66\textwidth}}
\toprule
\textbf{Term} & \textbf{Meaning in this lesson} \\
\midrule
OCR & Optical Character Recognition --- converts an image of printed or written text into editable text. \\
OMR & Optical Mark Recognition --- detects the presence or absence of marks at fixed positions on a form. \\
MICR & Magnetic Ink Character Recognition --- reads characters printed in magnetic ink, used on bank cheques. \\
Barcode & A one-dimensional pattern of parallel lines read optically to identify an item. \\
QR code & A two-dimensional matrix of squares (Quick Response) that holds more data than a barcode. \\
RFID & Radio Frequency Identification --- a tag and reader exchange radio waves; no line of sight is needed. \\
Biometrics & Input that identifies a person from a physical or behavioural trait (fingerprint, iris, face). \\
\bottomrule
\end{tabular}
\end{center}

\section{OCR --- Optical Character Recognition}
\textbf{OCR} stands for \textbf{Optical Character Recognition}. It takes a
scanned image of \emph{printed or handwritten text} and converts it into
machine-encoded, \emph{editable text}. The key word is
\emph{character}: OCR recognises individual letters, digits, and symbols, so
the output is text that a computer can search, copy, and edit. This is what
makes OCR different from OMR, which recognises only whether a position has
been marked and never identifies a letter.

OCR is used to digitise the printed word. A book, a bundle of invoices, a
passport page, or an old form can be scanned into an image, and OCR software
then extracts the text from that image. A useful way to remember the pairing
is: the \textbf{scanner} is the hardware that creates the image, while
\textbf{OCR} is the technology that turns the image of text into actual,
selectable text.

\begin{definitionbox}{OCR}
A recognition technology that reads characters from a scanned image of
printed or handwritten text and converts them into machine-encoded,
editable text.
\end{definitionbox}

OCR works best on clean, clearly printed text in a known font; its accuracy
drops when the page is skewed, smudged, or handwritten. When the input is
handwritten rather than printed, the same idea is sometimes called
\emph{Intelligent Character Recognition} (ICR), but for exam purposes both
count as OCR-style reading of text. OCR is \emph{not} a storage device and
\emph{not} a printer: it only converts an existing image into text.

\begin{example}
A student scans a page of a textbook into the computer. The scanner produces
an \emph{image}, so the words on the page cannot yet be copied or searched.
Running OCR on that image produces a text file in which every word can be
selected and edited. The scanner supplied the image; OCR supplied the text.
\end{example}

\section{OMR --- Optical Mark Recognition}
\textbf{OMR} stands for \textbf{Optical Mark Recognition}. It reads a form
by detecting the \emph{presence or absence of a mark} at fixed positions.
Typically the candidate fills a small bubble or shades a box with a pencil or
pen, and the OMR reader checks which positions have been darkened. Crucially,
OMR reads \emph{marks}, not characters: it does not know what a letter
``says'', only that a particular position has been filled.

Because it detects marks at fixed coordinates, OMR is ideal for forms where
the answer is one of a small set of choices. Its classic use is the
\textbf{MCQ answer sheet}: the machine reads which bubble has been shaded for
each question. It is also used for questionnaires, surveys, census forms, and
ballots. A very common exam trap is to confuse OMR with OCR; remember that
OMR works on \emph{marks}, while OCR works on \emph{text}.

\begin{definitionbox}{OMR}
A recognition technology that detects the presence or absence of marks at
fixed positions on a form, used mainly to read objective (MCQ) answer
sheets and survey forms.
\end{definitionbox}

OMR is extremely fast because it does not have to interpret a shape as a
character; it only has to decide whether a small region is dark or light.
That speed is why OMR is the standard method for large-scale examinations
where millions of answer sheets must be processed quickly. Its natural limit
is that it can record only \emph{which} option was marked, not any written
answer: if the candidate writes a sentence instead of shading a bubble, OMR
cannot read it.

\section{MICR --- Magnetic Ink Character Recognition}
\textbf{MICR} stands for \textbf{Magnetic Ink Character Recognition}. It
reads characters printed in a special \emph{magnetic ink} that contains iron
oxide. The characters are printed in a standard font along the bottom of a
\emph{bank cheque} as the \textbf{MICR line}, which normally encodes the
cheque number, the bank and branch code, and the account number. A MICR
reader magnetises the ink and detects the shapes of the characters as they
pass the read head.

Banking is the defining application of MICR. Cheque processing must be very
fast and highly secure, and MICR satisfies both needs: the reading is quick,
and because the ink is magnetic it is difficult to forge or to alter without
detection. When a question mentions a \emph{bank cheque} or \emph{magnetic
ink}, the answer is almost always MICR.

\begin{definitionbox}{MICR}
A recognition technology that reads characters printed in magnetic ink,
used mainly in banking to process the MICR line printed at the bottom of a
cheque.
\end{definitionbox}

Unlike ordinary ink, magnetic ink can be read even if the cheque is slightly
crumpled or over-stamped, because the reader senses the magnetic field of
the ink rather than a purely visual image. This ruggedness and the difficulty
of counterfeiting magnetic ink are exactly the reasons banks adopted MICR
for cheque clearing.

\section{The crucial trio: OCR vs OMR vs MICR}
These three names are the heart of the lesson, and JKSSB questions routinely
swap their meanings in the options. The table below is the single most useful
thing to memorise.

\begin{center}
\begin{tabular}{p{0.12\textwidth}p{0.26\textwidth}p{0.50\textwidth}}
\toprule
\textbf{Method} & \textbf{What it reads} & \textbf{Where it is used} \\
\midrule
OCR & Printed or written characters & Digitising books, documents, and passports \\
OMR & Marks on a form & MCQ answer sheets, surveys, ballots \\
MICR & Characters in magnetic ink & Bank cheques and cheque clearing \\
\bottomrule
\end{tabular}
\end{center}

The distinction to hold is: \textbf{MICR} uses \emph{magnetic ink};
\textbf{OCR} reads \emph{printed text}; \textbf{OMR} reads \emph{marks}. This
is the JKSSB trap TRAP-25, and the real paper item behind it (PYQ-30) asks
simply that magnetic ink means MICR.

It also helps to picture \emph{how} each one captures its data. OMR and OCR
both read an optical (light) signal: the reader sees the paper. MICR is the
odd one out, because its reader senses a \emph{magnetic} signal from the ink.
So a quick way to separate them in a hurry is: two are \emph{optical}, one is
\emph{magnetic}.

\begin{center}
\begin{tabular}{p{0.11\textwidth}p{0.20\textwidth}p{0.25\textwidth}p{0.31\textwidth}}
\toprule
\textbf{Method} & \textbf{Signal read} & \textbf{Input pattern} & \textbf{Cannot do} \\
\midrule
OCR & Optical (light) & Text characters & Cannot read marks alone \\
OMR & Optical (light) & Marks / shaded bubbles & Cannot read written characters \\
MICR & Magnetic & Magnetic-ink characters & Not used for forms or plain text \\
\bottomrule
\end{tabular}
\end{center}

\section{Barcodes and QR codes}
A \textbf{barcode} is a one-dimensional (1D) pattern of parallel lines of
varying width and spacing. A barcode scanner shines a light beam across the
lines and reads the reflected pattern, which represents a number or code
that identifies a product. Because the scanner must see the lines, a barcode
requires \emph{line of sight}. Barcodes are used on product labels, on the
ISBN of books, and on postal items.

A \textbf{QR code} --- \textbf{Quick Response} code --- is a
\emph{two-dimensional} (2D) matrix of black and white squares. It stores
more data than a barcode and can be read from any direction. Because it can
hold a URL, scanning a QR code can open a website or an app directly. Both a
barcode and a QR code are read \emph{optically}, so the nearest-neighbour
trap is simply that a QR code is two-dimensional and carries more
information, while a barcode is one-dimensional.

\begin{definitionbox}{QR code}
A two-dimensional matrix of black and white squares (Quick Response code)
that stores more data than a barcode and can be read from any direction.
\end{definitionbox}

\begin{example}
A shop assistant passes a product over a laser at the billing counter: the
scanner reads the 1D \emph{barcode} on the packet and looks up its price. A
customer points a phone camera at the square \emph{QR code} on a poster: the
phone reads the 2D code and opens the shop's website. The barcode carried a
short number; the QR code carried a whole web address.
\end{example}

\section{RFID --- Radio Frequency Identification}
\textbf{RFID} stands for \textbf{Radio Frequency Identification}. An RFID
system has two parts: a \emph{tag} attached to the object and a \emph{reader}
that interrogates it. The two exchange \emph{radio waves}, so no line of
sight is required, unlike a barcode. A passive tag carries no battery and is
powered by the reader's signal; an active tag has its own battery and can
broadcast over a longer range.

RFID is used where the tag may be out of sight or where many tags must be
read quickly: toll collection (for example FASTag), access-control cards,
passports, library books, and inventory tracking. The exam contrast is
always barcode/QR versus RFID: barcode and QR are read \emph{optically} and
need line of sight, while RFID uses \emph{radio waves} and can read a tag
without it being visible.

\section{Biometric input}
\textbf{Biometrics} identifies a person from a physical or behavioural
trait. The common physical traits are the \textbf{fingerprint}, the
\textbf{iris}, the \textbf{face}, and the palm; the voice is a common
behavioural trait. A sensor captures the trait and feeds the measured data to
the computer as \emph{input}; the system then compares it with the stored
template. Biometric input is a \emph{direct} form of input because the sensor
reads the trait itself rather than asking the user to type or encode it.

Biometrics are used for attendance systems, device unlocking, and
authentication, and they appear in national identity and benefit systems. A
useful checklist item for the exam is that a fingerprint reader, an iris
scanner, and a face-recognition camera are all \emph{input} devices; a device
that merely prints or displays the result is \emph{output}. Do not confuse
biometric traits (fingerprint, iris, face) with the optical reading machines:
a fingerprint sensor reads a part of the human body, not a paper code.

\section{Full-form ladder and memory hooks}
Direct-recall items on this topic usually ask for a full form, so learn the
expansions exactly. The same expansions appear on the slides; repeat them
until they can be produced without hesitation.

\begin{center}
\begin{tabular}{p{0.16\textwidth}p{0.72\textwidth}}
\toprule
\textbf{Short form} & \textbf{Full form} \\
\midrule
OCR & Optical Character Recognition \\
OMR & Optical Mark Recognition \\
MICR & Magnetic Ink Character Recognition \\
QR & Quick Response \\
RFID & Radio Frequency Identification \\
\bottomrule
\end{tabular}
\end{center}

Three one-line hooks carry most of the content of this lesson. First,
\emph{OCR reads text, OMR reads marks, MICR reads magnetic ink}. Second,
\emph{a barcode is one-dimensional, a QR code is two-dimensional, and RFID
uses radio waves}. Third, \emph{a biometric sensor reads a human trait ---
fingerprint, iris, face, or voice}. If a question seems unfamiliar, reduce
it to one of these three sentences and the answer usually becomes clear.

\section{Applications at a glance}
The table gathers the reading machines and code devices with their defining
application, for quick revision.

\begin{center}
\begin{tabular}{p{0.20\textwidth}p{0.66\textwidth}}
\toprule
\textbf{Technology} & \textbf{Defining application} \\
\midrule
OCR & Converting scanned printed text into editable text \\
OMR & Reading marks on objective (MCQ) answer sheets \\
MICR & Reading the magnetic-ink line at the bottom of a bank cheque \\
Barcode & Reading 1D product labels and book ISBNs \\
QR code & Reading 2D codes that can open a link or store more data \\
RFID & Reading tags by radio waves (tolls, access cards, inventory) \\
Biometrics & Identifying a person from a fingerprint, iris, face, or voice \\
\bottomrule
\end{tabular}
\end{center}

\section{Exam retrieval and common confusions}
Six distinctions prevent most errors on this topic.
\begin{enumerate}
  \item \textbf{MICR} uses \textbf{magnetic ink}; \textbf{OCR} reads
    \textbf{printed text}; \textbf{OMR} reads \textbf{marks}. Do not
    interchange them (TRAP-25, PYQ-30).
  \item A real PYQ asks which technology uses magnetic ink on cheques: the
    answer is \textbf{MICR} (PYQ-30).
  \item \textbf{OMR}, not OCR, is used for \textbf{MCQ answer sheets},
    because it reads marks at fixed positions.
  \item A \textbf{QR code} is two-dimensional and holds more data than a
    one-dimensional \textbf{barcode}; both are read optically.
  \item \textbf{RFID} uses \textbf{radio waves} and needs no line of sight;
    a barcode or QR code must be \emph{seen} by the reader.
  \item A \textbf{fingerprint} or \textbf{iris} scanner is \textbf{input};
    biometrics identify a person from a physical or behavioural trait.
\end{enumerate}

\begin{example}
An item states that a bank reads the characters at the bottom of a cheque
using OCR. This is wrong. The cheque characters are printed in magnetic ink,
so the technology is \textbf{MICR}. OCR would be used to read ordinary
\emph{printed text} from a scanned document, and OMR would be used to read
\emph{marks} on a form --- neither of which matches a cheque's magnetic-ink
line.
\end{example}

\begin{example}
An item lists four reading methods and asks which two are \emph{optical}.
OCR and OMR both use a light signal to read the paper, so they are the two
optical methods. MICR is \emph{magnetic}, not optical, and RFID is neither
optical nor magnetic --- it uses radio waves. Separating the methods by their
signal (optical, magnetic, radio) is a fast way to answer such a question.
\end{example}

\subsection*{Exercises}
\begin{enumerate}[label=\arabic*.]
  \item Expand OCR, OMR, and MICR, and state in one line what each reads.
  \item Which technology would read an objective (MCQ) answer sheet, and
    why?
  \item Explain why a bank uses MICR rather than OCR for cheque processing.
  \item Distinguish a barcode from a QR code on two points.
  \item State one way in which RFID differs from a barcode.
  \item Give three examples of biometric traits and state the purpose for
    which biometric input is commonly used.
  \item Classify these as optical, magnetic, or radio reading: OMR, MICR,
    RFID, OCR.
\end{enumerate}

\subsection*{Solutions}
\begingroup\small
\begin{enumerate}[label=\arabic*.]
  \item OCR = Optical Character Recognition, which reads printed or written
    characters; OMR = Optical Mark Recognition, which reads marks on a form;
    MICR = Magnetic Ink Character Recognition, which reads characters in
    magnetic ink.
  \item OMR, because it detects the presence or absence of a mark at fixed
    positions, which is exactly how a candidate fills bubbles on an MCQ
    answer sheet.
  \item The cheque characters are printed in magnetic ink, which MICR reads
    directly; the reading is fast, and the magnetic ink is difficult to
    forge or alter. OCR reads ordinary printed text, not magnetic-ink
    characters.
  \item A barcode is a one-dimensional pattern of parallel lines and holds
    less data; a QR code is a two-dimensional matrix of squares, holds more
    data, and can be read from any direction.
  \item RFID exchanges radio waves between a tag and a reader and needs no
    line of sight, whereas a barcode is read optically and must be visible
    to the scanner. RFID can also read many tags at once.
  \item Fingerprint, iris, and face (voice and palm are also valid).
    Biometric input is commonly used for attendance, authentication, and
    device unlocking.
  \item OMR and OCR are \emph{optical}; MICR is \emph{magnetic}; RFID is
    \emph{radio}.
\end{enumerate}
\endgroup
\end{document}
